\documentclass[11pt,a4paper,landscape]{article}
\usepackage[T1]{fontenc}
\usepackage[utf8]{inputenc}
\usepackage[french]{babel}
\usepackage{lmodern,amsmath,amssymb,graphicx,booktabs,xcolor}
\usepackage[margin=16mm]{geometry}
\usepackage{hyperref,fancyhdr}
\hypersetup{colorlinks=true,urlcolor=blue,linkcolor=blue,pdftitle={Plaque trouée : diagnostic des remappages successifs}}
\graphicspath{{../results/hudge_remap_diagnostic_20261115/figures/}}
\pagestyle{fancy}\fancyhf{}
\fancyhead[L]{WGPU-Matching — Diagnostic visuel de l'initialisation}
\fancyhead[R]{15 novembre 2026}\fancyfoot[R]{\thepage}
\setlength{\headheight}{14pt}
\setlength{\parindent}{0pt}\setlength{\parskip}{6pt}
\newcommand{\fullfig}[2]{\begin{center}\includegraphics[width=\linewidth,height=#2,keepaspectratio]{#1}\end{center}}
\begin{document}
{\LARGE\bfseries Plaque trouée : voir ce que le remappage corrige,\par
et ce qui empêche la propagation de continuer\par}
\vspace{5mm}
{\large Diagnostic sur trois paires indépendantes : niveaux 50, 150 et 191}

\textbf{Constat principal.} Une fenêtre non retenue n'est pas nécessairement une fenêtre
que le solveur n'arrive pas à recaler. Les images et les traces révèlent trois situations :
un recalage accepté par le solveur mais refusé par la cohérence avec le voisin,
un essai qui ne converge pas, et une fenêtre que la frontière de propagation n'atteint pas.
Les cartes de couverture seules mélangeaient ces situations.

\begin{center}
\begin{tabular}{lrrr}
\toprule
 & Niveau 50 & Niveau 150 & Niveau 191\\\midrule
Fenêtres éligibles (supports complets de $32\times32$ pixels) &36&36&36\\
Retenues avant la reprise affine &26&21&18\\
Retenues après la reprise et la propagation &34&27&18\\
Restantes non retenues &2&9&18\\
Parmi celles-ci, au moins un essai natif accepté par le solveur &2&8&13\\
Dernier rejet de propagation explicitement lié au gradient &2&7&7\\
\bottomrule
\end{tabular}
\end{center}

\textbf{Exemple décisif : fenêtre 102, niveau 191.} Le dernier essai fait passer la
corrélation de 0,713 à 0,889. Le solveur accepte sa sortie. La correction du déplacement
est de 4,96 pixels, inférieure au seuil de 8 pixels, mais la correction du gradient vaut
0,403, supérieure au seuil de 0,25 : la propagation ne la retient pas.
Les agrandissements de la page~\pageref{page:191patch} montrent le mouchetis correspondant.

\textbf{Portée.} Il s'agit d'un diagnostic d'initialisation de la corrélation d'images
numériques (DIC), avec le solveur Gauss--Newton à composition inverse (IC-GN).
Une bonne corrélation et un statut accepté ne prouvent pas l'exactitude mécanique.
Aucun champ de déplacement de référence associé à ces images n'a été identifié ;
le fichier de déplacements de l'éprouvette en croix n'est pas utilisé.

\newpage
\section*{1. Ce qui a été remappé et enregistré}
Chaque cible est comparée directement à \texttt{Img\_F.tif}. Les niveaux 50, 150 et 191
sont trois expériences séparées, sans exploitation des images intermédiaires.
Les fichiers source sont les images \texttt{Img\_wa\_g050.tif},
\texttt{Img\_wa\_g150.tif} et \texttt{Img\_wa\_g191.tif} du dossier
\texttt{HudgeStrains/FEA}. Ce sont des indices de fichier, pas des pourcentages de traction.

À chaque tour $k$, une transformation globale $A_k\mathbf x+\mathbf t_k$ est ajustée
sur les centres déjà retenus. La cible originale $I_1$ est ramenée dans le repère de
la référence $I_0$ :
\[
 J_k(\mathbf z)=I_1(A_k\mathbf z+\mathbf t_k),\qquad
 \boldsymbol\varphi(\mathbf x)=A_k\bigl(\mathbf x+\mathbf r_k(\mathbf x)\bigr)+\mathbf t_k,
 \qquad F=A_k(I+\nabla\mathbf r_k).
\]
Le solveur recherche un résidu local $\mathbf r_k$ dans l'image rectifiée. Une sortie
acceptable est recomposée, puis résolue à nouveau sur les pixels de la cible originale.
La confirmation des graines et la propagation appliquent ensuite leurs critères existants.
\textbf{Chaque image rectifiée repart de la cible originale}, pas de l'image rectifiée
du tour précédent : on ne cumule pas les interpolations d'image entre les tours.

\begin{minipage}[t]{.48\linewidth}
\textbf{Paramètres inchangés}\par
Supports et pas de 32 pixels ; interpolation bilinéaire ; 30 itérations au maximum ;
tolérance $10^{-3}$ ; corrélation normalisée centrée minimale de 0,65.
La cohérence en propagation exige
\[
 \|\mathbf d_{\text{sortie}}-\mathbf d_{\text{proposée}}\|_2\leq8\ {\rm pixels},\quad
 \|G_{\text{sortie}}-G_{\text{proposé}}\|_F\leq0,25.
\]
Ici $G=F-I$ et $\|\cdot\|_F$ est la norme de Frobenius.
Ces critères ne s'appliquent pas indistinctement à toutes les étapes : la sélection
initiale des hypothèses utilise le statut du solveur.
\end{minipage}\hfill
\begin{minipage}[t]{.48\linewidth}
\textbf{Lecture des figures}\par
Les vues globales montrent le même rectangle de référence, $440\times225$ pixels.
Les fichiers remappés complets restent en $1024\times1024$ pixels.
Les agrandissements utilisent les centres exacts et les 32 pixels du support.
Les gris sont fixes entre 0 et 255 ; les résidus absolus sont fixes entre 0 et 128,
saturés au-delà. Les résidus ne sont pas corrigés photométriquement.
La corrélation affichée est recalculée après l'essai, sur les pixels originaux ;
elle n'est pas une lecture de la dernière valeur interne du solveur.
\end{minipage}

Les relances instrumentées reproduisent \textbf{exactement les masques retenus et tous
les paramètres numériques archivés} pour les trois niveaux. Aucun seuil n'a été modifié.
Les figures affichent aussi les propositions non validées, explicitement identifiées.

\newpage
\section*{2. Niveau 50 : la reprise récupère huit fenêtres}
\fullfig{sequence_50.pdf}{72mm}
\fullfig{map_50.pdf}{57mm}
Le remappage global ne doit pas être confondu avec la solution locale : même si les
contours sont proches, le mouchetis peut rester déphasé. Les fenêtres vertes sont les huit
récupérations ; les deux oranges persistent après la reprise. Les nombres sont les
identifiants de fenêtres, conservés dans les exports détaillés.

\newpage
\section*{3. Niveau 50 : une récupération et un rejet malgré un bon recalage}
\fullfig{patch_50_35.pdf}{57mm}
\textbf{Fenêtre 35, récupérée.} La proposition affine globale seule est mauvaise localement
(corrélation $-0,074$). Une proposition transportée par le voisinage permet au solveur
d'atteindre 0,953 ; cette fenêtre est finalement retenue. L'ajustement global n'explique
donc pas à lui seul la récupération.
\fullfig{patch_50_36.pdf}{57mm}
\textbf{Fenêtre 36, non retenue.} Le recalage final présente un mouchetis beaucoup plus proche
(0,937), mais sa correction du gradient vaut 0,433. Le dernier rejet est celui du seuil
de cohérence 0,25, pas celui du seuil photométrique 0,65. Une similarité élevée ne suffit
cependant pas à autoriser automatiquement cette sortie.

\newpage
\section*{4. Niveau 150 : le remappage global ne résout pas toute l'hétérogénéité}
\fullfig{sequence_150.pdf}{72mm}
\fullfig{map_150.pdf}{57mm}
Le résultat passe de 21 à 27 fenêtres retenues. Les résidus visibles autour des régions
locales montrent ce que l'affine globale laisse au solveur. Parmi les neuf fenêtres
restantes, huit possèdent au moins un essai accepté par le solveur natif. Le dernier
état de la propagation indique sept rejets de gradient, un rejet du solveur et une
fenêtre sans motif de rejet enregistré à cette étape.

\newpage
\section*{5. Niveau 150 : séparer l'efficacité du solveur et la décision de propagation}
\fullfig{patch_150_35.pdf}{57mm}
\textbf{Fenêtre 35, récupérée.} La dernière proposition locale est déjà proche (0,811).
Le solveur améliore encore le mouchetis (0,899), avec une correction du gradient 0,183.
\fullfig{patch_150_102.pdf}{57mm}
\textbf{Fenêtre 102, non retenue.} L'essai natif améliore la corrélation de 0,720 à 0,900.
La correction du déplacement reste modérée (4,10 pixels), mais celle du gradient
atteint 0,332 : le contrôle de cohérence arrête cette branche.
Les résidus montrent qu'il existe ici une sortie photométriquement crédible à examiner,
plutôt qu'une absence manifeste de motif correspondant.

\newpage
\section*{6. Niveau 150 : ce que nous ne pouvons pas appeler un échec du solveur}
\fullfig{patch_150_164.pdf}{64mm}
\textbf{Fenêtre 164, centre $(143,5;175,5)$.} Aucune tentative n'est enregistrée dans cette
reprise pour cette fenêtre. Les deux panneaux montrent uniquement les propositions
initiales ; ils ne montrent pas une sortie du solveur. La frontière ne l'atteint pas.
Sa corrélation initiale faible indique que cette proposition ne suffit pas, mais ne
démontre pas qu'une initialisation voisine adéquate échouerait.
\fullfig{patch_150_70.pdf}{57mm}
\textbf{Contrôle retenu : fenêtre 70.} La texture est correctement rapprochée et la corrélation
atteint 0,934. Cette fenêtre était déjà retenue avant la reprise : son état est conservé.
Un nouvel essai peut avoir une correction du gradient supérieure au seuil sans annuler
une acceptation antérieure. On ne peut donc pas déduire le statut final du seul dernier essai.

\newpage
\section*{7. Niveau 191 : pourquoi répéter la même affine s'arrête}
\fullfig{sequence_191.pdf}{72mm}
\fullfig{map_191.pdf}{57mm}
Le premier tour ne crée aucune fenêtre supplémentaire : l'arrêt intervient à 18/36.
L'affine ajustée aux 18 centres retenus possède notamment $A_{22}=4,871$ ; le résidu
médian d'ajustement des centres est de 4,18 pixels. C'est une approximation globale
des observations retenues, pas une mesure locale uniforme de déformation.
Réutiliser exactement le même ensemble et la même règle ne fournit pas de nouvelle
information pour le tour suivant. Le diagnostic doit porter sur les branches bloquées.

\newpage\label{page:191patch}
\section*{8. Niveau 191 : le même motif est retrouvé, mais la branche est refusée}
\fullfig{patch_191_102.pdf}{64mm}
\textbf{Fenêtre 102.} Le mouchetis après résolution est visiblement plus proche de la référence.
Le statut natif est accepté ; la corrélation vaut 0,889. Le dernier rejet de propagation
provient de $\Delta G=0,403>0,25$. Ce constat identifie le verrou logiciel de cet essai,
sans établir que la sortie constitue une mesure mécanique exacte.
\fullfig{patch_191_36.pdf}{57mm}
\textbf{Fenêtre 36.} Aucune tentative de reprise n'est enregistrée à cet endroit.
Une région non atteinte ne doit pas être comptée comme un échec photométrique démontré.
Il faut d'abord apporter une proposition locale qui permette réellement de la tester.

\newpage
\section*{9. Conclusions opérationnelles et limites}
\textbf{Ce que ces images établissent.} Le remappage global rapproche la géométrie,
mais ne garantit pas la coïncidence des mouchetis localement. Les propositions issues
des voisins peuvent permettre à IC-GN de retrouver une texture très proche.
Le mécanisme de cohérence arrête ensuite certaines sorties acceptées par le solveur.
Les régions situées derrière ces arrêts peuvent ne jamais être essayées.

\textbf{Ce que ces images n'établissent pas.} On ne dispose pas de vérité de déplacement
pour ces paires. Une corrélation élevée ne garantit ni les six paramètres affines exacts,
ni l'absence d'une erreur localisée. Le résidu seul ne permet pas de départager toutes
les causes : erreur de géométrie, interpolation, changement de contraste, bruit ou
inadéquation du modèle affine sur le support. Le présent diagnostic ne mesure pas
la part de chacune.

\textbf{Prochaine expérience justifiée.} Reprendre les sorties photométriquement crédibles
mais rejetées par le gradient, telles que la fenêtre 102, comme \emph{hypothèses à vérifier}
avec une confirmation indépendante et des supports adaptés. Mesurer si elles ouvrent
la frontière sans introduire de faux résultats sur des cas à vérité connue. Il ne faut
pas simplement supprimer le seuil de cohérence pour augmenter la couverture.
Une nouvelle affine globale seule ne cible pas ce verrou.

\textbf{Traçabilité.} Le dossier
\texttt{results/hudge\_remap\_diagnostic\_20261115/} contient les remappages complets,
les figures, les exports par fenêtre et les appels enregistrés dans
\texttt{recorded/rectified\_native/trace\_*/}. Chaque appel fournit les indices,
paramètres proposés et résolus, statut, étape et numéro de tour ; les appels directs
du solveur contiennent aussi les nombres d'itérations.
Les codes de dernier rejet de propagation sont : 0, aucun motif enregistré ;
1, solveur non accepté ou sortie non finie ; 2, correction du déplacement ;
3, correction du gradient. Un code nul n'est pas une preuve de réussite.

\textbf{Reproduction.} Les commandes exactes, les empreintes des entrées et les contrôles
de parité figurent dans le fichier \texttt{README.md} et le manifeste de ce dossier.
Le script \texttt{diagnose\_hudge\_remaps.py} ne relance pas l'algorithme et n'utilise
aucune vérité analytique : il reconstruit les vues depuis les traces enregistrées.
Les exemples locaux sont choisis par catégorie puis proximité du centre fixe $(175;110)$,
sans recherche du meilleur score. Les traces complètes couvrent les 36 fenêtres de chaque paire.

\end{document}
